\frfilename{mt2a4.tex}
\versiondate{4.3.14}
\copyrightdate{1996}

\def\chaptername{Lampiran}
\def\sectionname{Ruang bernorma}

\newsection{2A4}

Dalam Bab 24 saya membahas ruang-ruang $L^p$, untuk
$1\le p\le\infty$, dan menguraikan sifat-sifatnya yang paling mendasar.
Ruang-ruang ini membentuk sekelompok contoh fundamental bagi teori umum
ruang bernorma, yang merupakan landasan analisis fungsional. Buku ini
bukanlah buku yang seharusnya Anda gunakan untuk mempelajari teori
tersebut, tetapi sekali lagi, mungkin akan memudahkan Anda jika saya
menguraikan secara ringkas bagian-bagian teori umum yang mutlak diperlukan
untuk memahami gagasan-gagasan di sini.

\leader{2A4A}{Lapangan real dan kompleks} Meskipun bagian-bagian
terpenting teori ini, dari sudut pandang teori ukur, paling efektif
ditangani dalam ruang linear {\it real}, terdapat banyak penerapan yang
mutlak memerlukan ruang linear {\it kompleks}. Karena itu, saya akan
menggunakan ungkapan

\Centerline{“$U$ adalah ruang linear atas $\RoverC$”}

\noindent untuk menyatakan bahwa $U$ adalah ruang linear atas lapangan
$\Bbb R$ atau ruang linear atas lapangan $\Bbb C$; dipahami bahwa dalam
setiap konteks tertentu, semua ruang linear yang dibahas berada di atas
lapangan yang sama. Dengan cara serupa, saya akan menulis
“$\alpha\in\RoverC$” untuk menyatakan bahwa $\alpha$ termasuk dalam
lapangan dasar yang sedang digunakan.

\leader{2A4B}{Definisi (a)} Sebuah {\bf ruang bernorma} adalah ruang
linear $U$ atas $\RoverC$ yang dilengkapi dengan sebuah {\bf norma},
yaitu fungsional $\|\,\|:U\to\coint{0,\infty}$ sedemikian sehingga

\qquad $\|u+v\|\le\|u\|+\|v\|$ untuk semua $u$, $v\in U$,

\qquad $\|\alpha u\|=|\alpha|\|u\|$ untuk $u\in U$,
$\alpha\in\RoverC$,

\qquad $\|u\|=0$ hanya jika $u=0$, yaitu vektor nol di $U$.

\cmmnt{\noindent (Perhatikan bahwa jika $u=0$ (vektor nol), maka
$0u=u$ (dengan $0$ yang ini adalah skalar nol), sehingga
$\|u\|=|0|\|u\|=0$.)
}

\header{2A4Bb}{\bf (b)} Jika $U$ adalah ruang bernorma, maka kita
memperoleh metrik $\rho$ pada $U$ dengan menetapkan bahwa
$\rho(u,v)=\|u-v\|$ untuk $u$, $v\in U$. \prooflet{\Prf\
$\rho(u,v)\in\coint{0,\infty}$ untuk semua $u$, $v$, karena
$\|u\|\in\coint{0,\infty}$ untuk setiap $u$.
$\rho(u,v)=\rho(v,u)$ untuk semua $u$, $v$, karena
$\|v-u\|=|-1|\|u-v\|=\|u-v\|$ untuk semua $u$, $v$.
Jika $u$, $v$, $w\in U$, maka

\Centerline{$\rho(u,w)=\|u-w\|=\|(u-v)+(v-w)\|\le\|u-v\|+\|v-w\|
=\rho(u,v)+\rho(v,w)$.}

\noindent Jika $\rho(u,v)=0$, maka $\|u-v\|=0$, sehingga $u-v=0$ dan
$u=v$.\ \Qed}

\cmmnt{Dengan demikian, kita memperoleh topologi yang bersesuaian,
beserta himpunan terbuka dan tertutup, penutupan, barisan konvergen, dan
sebagainya.}

\spheader 2A4Bc Jika $U$ adalah ruang bernorma, suatu himpunan
$A\subseteq U$ disebut {\bf terbatas}\cmmnt{ (terhadap normanya)} jika
$\{\|u\|:u\in A\}$ terbatas di $\Bbb R$\cmmnt{; yaitu, terdapat
$M\ge 0$ sehingga $\|u\|\le M$ untuk setiap $u\in A$}.

\leader{2A4C}{Subruang linear (a)} Jika $U$ adalah sebarang ruang
bernorma dan $V$ adalah subruang linear dari $U$, maka $V$ juga
merupakan ruang bernorma, apabila norma pada $V$ diambil sebagai
pembatasan pada $V$ dari norma $U$\cmmnt{; verifikasinya langsung}.

\spheader 2A4Cb Jika $V$ adalah subruang linear dari $U$, maka
penutupannya\cmmnt{ $\overline{V}$} juga demikian.
\prooflet{\Prf\ Ambil $u$, $u'\in\overline{V}$ dan
$\alpha\in\RoverC$. Jika $\epsilon>0$, tetapkan
$\delta=\epsilon/(2+|\alpha|)>0$; maka terdapat $v$, $v'\in V$
sedemikian sehingga $\|u-v\|\le\delta$ dan
$\|u'-v'\|\le\delta$. Sekarang $v+v'$, $\alpha v\in V$ dan

\Centerline{$\|(u+u')-(v+v')\|\le\|u-v\|+\|u'-v'\|\le\epsilon$,
\quad$\|\alpha u-\alpha v\|\le|\alpha|\|u-v\|\le\epsilon$.}

\noindent Karena $\epsilon$ sebarang, $u+u'$ dan $\alpha u$ termasuk
dalam $\overline{V}$; karena $u$, $u'$, dan $\alpha$ sebarang, dan
$0$ tentu termasuk dalam $V\subseteq\overline{V}$, maka
$\overline{V}$ adalah subruang linear dari $U$.\ \Qed}

\leader{2A4D}{Ruang Banach (a)} Jika $U$ adalah ruang bernorma, suatu
barisan $\sequencen{u_n}$ di $U$ disebut {\bf Cauchy} jika
$\|u_m-u_n\|\to 0$ ketika $m$, $n\to\infty$\cmmnt{, yaitu, untuk setiap
$\epsilon>0$ terdapat $n_0\in\Bbb N$ sedemikian sehingga
$\|u_m-u_n\|\le\epsilon$ untuk semua $m$, $n\ge n_0$}.

\spheader 2A4Db Sebuah ruang bernorma $U$ disebut {\bf lengkap} jika
setiap barisan Cauchy mempunyai limit; ruang bernorma lengkap disebut
{\bf ruang Banach}.

\leader{2A4E}{}\cmmnt{ Hasil berikut berguna untuk diketahui.

\medskip

\noindent}{\bf Lema} Misalkan $U$ adalah ruang bernorma sedemikian
sehingga $\sequencen{u_n}$ konvergen\cmmnt{ (yaitu, mempunyai limit)}
di $U$ setiap kali $\sequencen{u_n}$ adalah barisan di $U$ yang memenuhi
$\|u_{n+1}-u_n\|\le 4^{-n}$ untuk setiap $n\in\Bbb N$. Maka $U$
lengkap.

\proof{ Misalkan $\sequencen{u_n}$ adalah sebarang barisan Cauchy di
$U$. Untuk setiap $k\in\Bbb N$, ambil $n_k\in\Bbb N$ sedemikian
sehingga $\|u_{m}-u_n\|\le 4^{-k}$ apabila $m$, $n\ge n_k$. Tetapkan
$v_k=u_{n_k}$ untuk setiap $k$.
Maka $\|v_{k+1}-v_k\|\le 4^{-k}$ (baik $n_k\le n_{k+1}$ maupun
$n_{k+1}\le n_k$). Jadi $\sequence{k}{v_k}$ mempunyai limit $v\in U$.
Saya akan menunjukkan bahwa $v$ adalah limit yang diperlukan dari
$\sequencen{u_n}$. Diberikan $\epsilon>0$, ambil $l\in\Bbb N$
sedemikian sehingga $\|v_k-v\|\le\epsilon$ untuk setiap $k\ge l$;
ambil $k\ge l$ sedemikian sehingga $4^{-k}\le\epsilon$; maka jika
$n\ge n_k$,

\Centerline{$\|u_n-v\|=\|(u_n-v_k)+(v_k-v)\|
\le\|u_n-v_k\|+\|v_k-v\|
\le\|u_n-u_{n_k}\|+\epsilon
\le 2\epsilon$.}

\noindent Karena $\epsilon$ sebarang, $v$ adalah limit dari
$\sequencen{u_n}$. Karena $\sequencen{u_n}$ sebarang, $U$ lengkap.
}%end of proof of 2A4E

\leader{2A4F}{Operator linear terbatas (a)} Misalkan $U$, $V$ adalah
dua ruang bernorma. Sebuah operator linear $T:U\to V$ disebut
{\bf terbatas} jika $\{\|Tu\|:u\in U,\,\|u\|\le 1\}$ terbatas.
\cmmnt{({\bf Peringatan!} Dalam konteks ini, kita tidak mensyaratkan
agar seluruh himpunan nilai $T[U]$ terbatas; sebuah “operator linear
terbatas” tidak harus merupakan apa yang biasanya kita sebut “fungsi
terbatas”.)}
Tuliskan $\eurm B(U;V)$ untuk ruang semua operator linear terbatas dari
$U$ ke $V$, dan untuk $T\in \eurm B(U;V)$ tuliskan
$\|T\|= \sup\{\|Tu\|:u\in U,\,\|u\|\le 1\}$.

\header{2A4Fb}{\bf (b)}
\cmmnt{Fakta yang berguna:} $\|Tu\|\le\|T\|\|u\|$ apabila
$T\in\eurm B(U;V)$ dan $u\in U$. \prooflet{\Prf\ Jika
$|\alpha|>\|u\|$, maka

\Centerline{$\|\Bover1{\alpha}u\|=\Bover1{|\alpha|}\|u\|\le 1$,}

\noindent sehingga

\Centerline{$\|Tu\|=\|\alpha T(\Bover1{\alpha}u)\|
=|\alpha|\|T(\Bover1{\alpha}u)\|\le|\alpha|\|T\|$;}

\noindent karena $\alpha$ sebarang, $\|Tu\|\le\|T\|\|u\|$.\ \Qed}

\spheader 2A4Fc Sebuah operator linear $T:U\to V$ terbatas jika dan
hanya jika operator itu kontinu untuk topologi norma pada $U$ dan $V$.
\prooflet{\Prf\
(i) Jika $T$ terbatas, $u_0\in U$, dan $\epsilon>0$, maka

\Centerline{$\|Tu-Tu_0\|=\|T(u-u_0)\|\le\|T\|\|u-u_0\|\le\epsilon$}

\noindent apabila $\|u-u_0\|\le\Bover{\epsilon}{1+\|T\|}$; menurut
2A3H, $T$ kontinu. (ii) Jika $T$ kontinu, maka terdapat $\delta>0$
sedemikian sehingga $\|Tu\|=\|Tu-T0\|\le 1$ apabila
$\|u\|=\|u-0\|\le\delta$. Jika sekarang $\|u\|\le 1$,

\Centerline{$\|Tu\|=\Bover1{\delta}\|T(\delta u)\|\le\Bover1{\delta}$,}

\noindent sehingga $T$ adalah operator terbatas.\ \Qed}

\spheader 2A4Fd\dvAnew{2013} Jika $U$, $V$, dan $W$ adalah ruang
bernorma, $S\in\eurm B(U;V)$ dan $T\in\eurm B(V;W)$, maka
$TS\in\eurm B(U;W)$ dan $\|TS\|\le\|T\|\|S\|$.
\prooflet{\Prf\ Saya mengandaikan bahwa Anda sudah mengetahui, tetapi
bagaimanapun juga mudah diperiksa, bahwa $TS:U\to W$ adalah operator
linear. Sekarang, jika $u\in U$ dan $\|u\|\le 1$,

\Centerline{$\|TSu\|=\|T(Su)\|\le\|T\|\|Su\|\le\|T\|\|S\|$}

\noindent (menggunakan (b) untuk pertidaksamaan di tengah), sehingga
$TS$ terbatas dan $\|TS\|\le\|T\|\|S\|$.\ \Qed}

\leader{2A4G}{Teorema} $\eurm B(U;V)$ adalah ruang linear atas
$\RoverC$, dan $\|\,\|$ adalah norma pada $\eurm B(U;V)$.

\proof{ Seperti dalam 2A4Fd, mudah diperiksa bahwa jika $S:U\to V$ dan
$T:U\to V$ adalah operator linear, dan $\alpha\in\RoverC$, maka kita
memperoleh operator linear $S+T$ dan $\alpha T$ dari $U$ ke $V$ yang
didefinisikan oleh rumus

\Centerline{$(S+T)(u)=Su+Tu$,\quad$(\alpha T)(u)=\alpha(Tu)$}

\noindent untuk setiap $u\in U$; selain itu, menurut definisi
penjumlahan dan perkalian skalar ini, ruang semua operator linear dari
$U$ ke $V$ merupakan ruang linear. Sekarang terlihat bahwa apabila
$S$, $T\in\eurm B(U;V)$, $\alpha\in\RoverC$, $u\in U$, dan
$\|u\|\le 1$,

\Centerline{$\|(S+T)(u)\|=\|Su+Tu\|\le\|Su\|+\|Tu\|\le\|S\|+\|T\|$,}

\Centerline{$\|(\alpha
T)u\|=\|\alpha(Tu)\|=|\alpha|\|Tu\|\le|\alpha|\|T\|$;}

\noindent sehingga $S+T$ dan $\alpha T$ termasuk dalam
$\eurm B(U;V)$, dengan $\|S+T\|\le\|S\|+\|T\|$ dan
$\|\alpha T\|\le|\alpha|\|T\|$. Ini menunjukkan bahwa
$\eurm B(U;V)$ adalah subruang linear dari ruang semua operator linear,
dan karena itu merupakan ruang linear atas $\RoverC$. Untuk memeriksa
bahwa rumus yang diberikan untuk $\|T\|$ mendefinisikan norma, sebagian
besar pekerjaan telah selesai; demi kelengkapan formal, perlu saya
catat bahwa $\|T\|\in\coint{0,\infty}$ untuk setiap $T$; jika
$\alpha=0$, maka tentu saja $\|\alpha T\|=0=|\alpha|\|T\|$; untuk
$\alpha$ lainnya,

\Centerline{$|\alpha|\|T\|=|\alpha|\|\alpha^{-1}\alpha T\|
\le|\alpha||\alpha^{-1}|\|\alpha T\|=\|\alpha T\|\le|\alpha|\|T\|$,}

\noindent sehingga $\|\alpha T\|=|\alpha|\|T\|$. Akhirnya, jika
$\|T\|=0$, maka $\|Tu\|\le\|T\|\|u\|=0$ untuk setiap $u\in U$,
sehingga $Tu=0$ untuk setiap $u$ dan $T$ adalah operator nol (dalam
ruang semua operator linear, dan dengan demikian juga dalam subruangnya
$\eurm B(U;V)$).
}%end of proof of 2A4G

\leader{2A4H}{Ruang dual} Kasus terpenting dari $\eurm B(U;V)$ adalah
ketika $V$ merupakan lapangan skalar $\RoverC$ itu sendiri\cmmnt{
(tentu saja kita dapat memandang $\RoverC$ sebagai ruang bernorma atas
dirinya sendiri, dengan menulis $\|\alpha\|=|\alpha|$ untuk setiap
skalar $\alpha$)}. Dalam hal ini, $\eurm B(U;\RoverC)$ disebut
{\bf dual} dari $U$; ruang ini lazim dinotasikan dengan $U'$ atau
$U^*$; saya menggunakan yang terakhir.

\leader{2A4I}{Perluasan operator terbatas: Teorema} Misalkan $U$
adalah ruang bernorma dan $V\subseteq U$ subruang linear yang rapat.
Misalkan $W$ adalah ruang Banach dan $T_0:V\to W$ operator linear
terbatas; maka terdapat tepat satu operator linear terbatas $T:U\to W$
yang memperluas $T_0$, dan $\|T\|=\|T_0\|$.

\proof{{\bf (a)} Untuk setiap $u\in U$, terdapat barisan
$\sequencen{v_n}$ di $V$ yang konvergen ke $u$. Sekarang

\Centerline{$\|T_0v_m-T_0v_n\|=\|T_0(v_m-v_n)\|
\le\|T_0\|\|v_m-v_n\|\le\|T_0\|(\|v_m-u\|+\|u-v_n\|)
\to 0$}

\noindent ketika $m$, $n\to\infty$, sehingga
$\sequencen{T_0v_n}$ adalah barisan Cauchy dan
$w=\lim_{n\to\infty}T_0v_n$ terdefinisi di $W$. Jika
$\sequencen{v'_n}$ adalah barisan lain di $V$ yang konvergen ke $u$,
maka

$$\eqalign{\|w-T_0v'_n\|
&\le\|w-T_0v_n\|+\|T_0(v_n-v'_n)\|\cr
&\le\|w-T_0v_n\|+\|T_0\|(\|v_n-u\|+\|u-v'_n\|)\to 0\cr}$$

\noindent ketika $n\to\infty$, sehingga $w$ juga merupakan limit dari
$\sequencen{T_0v'_n}$.

\medskip

{\bf (b)} Karena itu, kita dapat mendefinisikan $T:U\to W$ dengan
menetapkan $Tu=\lim_{n\to\infty}T_0v_n$ apabila
$\sequencen{v_n}$ adalah barisan di $V$ yang konvergen ke $u$. Jika
$v\in V$, kita dapat menetapkan $v_n=v$ untuk setiap $n$ dan melihat
bahwa $Tv=T_0v$; jadi $T$ memperluas $T_0$. Jika $u$, $u'\in U$ dan
$\alpha\in\RoverC$, ambil barisan $\sequencen{v_n}$,
$\sequencen{v'_n}$ di $V$ yang masing-masing konvergen ke $u$, $u'$;
dalam hal ini

\Centerline{$\|(u+u')-(v_n+v'_n)\|\le\|u-v_n\|+\|u'-v'_n\|\to 0$,
\quad$\|\alpha u-\alpha v_n\|=|\alpha|\|u-v_n\|\to 0$}

\noindent ketika $n\to\infty$, sehingga
$T(u+u')=\lim_{n\to\infty}T_0(v_n+v'_n)$,
$T(\alpha u)=\lim_{n\to\infty}T_0(\alpha v_n)$, dan

$$\eqalign{\|T(u+u')-Tu-Tu'\|
&\le\|T(u+u')-T_0(v_n+v'_n)\|+\|T_0v_n-Tu\|+\|T_0v'_n-Tu'\|\cr
&\to 0,\cr}$$

\Centerline{$\|T(\alpha u)-\alpha Tu\|
\le\|T(\alpha u)-T_0(\alpha v_n)\|+|\alpha|\|T_0v_n-Tu\|\to 0$}

\noindent ketika $n\to\infty$. Ini berarti
$\|T(u+u')-Tu-Tu'\|=0$, $\|T(\alpha u)-\alpha Tu\|=0$, sehingga
$T(u+u')=Tu+Tu'$ dan $T(\alpha u)=\alpha Tu$; karena $u$, $u'$, dan
$\alpha$ sebarang, $T$ linear.

\medskip

{\bf (c)} Untuk setiap $u\in U$, misalkan $\sequencen{v_n}$ adalah
barisan di $V$ yang konvergen ke $u$. Maka

$$\eqalign{\|Tu\|
&\le\|T_0v_n\|+\|Tu-T_0v_n\|
\le\|T_0\|\|v_n\|+\|Tu-T_0v_n\|\cr
&\le\|T_0\|(\|u\|+\|v_n-u\|)+\|Tu-T_0v_n\|
\to\|T_0\|\|u\|\cr}$$

\noindent ketika $n\to\infty$, sehingga
$\|Tu\|\le\|T_0\|\|u\|$. Karena $u$ sebarang, $T$ terbatas dan
$\|T\|\le\|T_0\|$. Tentu saja $\|T\|\ge\|T_0\|$ karena $T$
memperluas $T_0$.

\medskip

{\bf (d)} Akhirnya, misalkan $\tilde T$ adalah operator linear terbatas
lain dari $U$ ke $W$ yang memperluas $T_0$. Jika $u\in U$, terdapat
barisan $\sequencen{v_n}$ di $V$ yang konvergen ke $u$; sekarang

\Centerline{$\|\tilde Tu-Tu\|\le\|\tilde T(u-v_n)\|+\|T(v_n-u)\|
\le(\|\tilde T\|+\|T\|)\|u-v_n\|\to 0$}

\noindent ketika $n\to\infty$, sehingga $\|\tilde Tu-Tu\|=0$ dan
$\tilde Tu=Tu$. Karena $u$ sebarang, $\tilde T=T$. Jadi $T$ tunggal.
}%end of proof of 2A4I

\leader{2A4J}{Aljabar bernorma (a)} Sebuah {\bf aljabar bernorma}
adalah ruang bernorma $(U,\|\,\|)$ yang dilengkapi dengan suatu
perkalian, yaitu operator biner $\times$ pada $U$, sedemikian sehingga

\Centerline{$u\times(v\times w)=(u\times v)\times w$,}

\Centerline{$u\times(v+w)=(u\times v)+(u\times w)$,
\quad$(u+v)\times w=(u\times w)+(v\times w)$,}

\Centerline{$(\alpha u)\times v=u\times(\alpha v)=\alpha(u\times v)$,}

\Centerline{$\|u\times v\|\le\|u\|\|v\|$}

\noindent untuk semua $u$, $v$, $w\in U$ dan $\alpha\in\RoverC$.

\spheader 2A4Jb Sebuah {\bf aljabar Banach} adalah aljabar bernorma
yang merupakan ruang Banach. Sebuah aljabar bernorma\cmmnt{ $U$}
disebut {\bf komutatif} jika perkaliannya komutatif\cmmnt{, yaitu,
$u\times v=v\times u$ untuk semua $u$, $v\in U$}.

\leader{*2A4K}{Definisi} Sebuah ruang bernorma $U$ disebut
{\bf konveks seragam} jika untuk setiap $\epsilon>0$ terdapat
$\delta>0$ sedemikian sehingga $\|u+v\|\le 2-\delta$ apabila
$u$, $v\in U$, $\|u\|=\|v\|=1$, dan $\|u-v\|\ge\epsilon$.

\discrpage
